\documentclass{article}

% Language setting
% Replace `english' with e.g. `spanish' to change the document language
\usepackage[english]{babel}
\usepackage{float}

% Set page size and margins
% Replace `letterpaper' with `a4paper' for UK/EU standard size
\usepackage[letterpaper,top=2cm,bottom=2cm,left=3cm,right=3cm,marginparwidth=1.75cm]{geometry}

% Useful packages
\usepackage{amsmath}
\usepackage{graphicx}
\usepackage[colorlinks=true, allcolors=blue]{hyperref}

\title{Labwork 5: Neuron Network with Backprobagation}
\author{Nguyen Nhu Duy - 2540044}

\begin{document}
\setlength{\parindent}{0pt}
\maketitle
\section{Backpropagation}

The algorithm computes the gradient of the loss function with respect to each weight and bias using the chain rule, then updates the parameters through gradient descent.

\subsection{Output Layer Error}

For the output layer, the error term (delta) is calculated as:

\[
\delta = \hat{y} - y
\]

This simplified form is obtained from combining the derivative of BCE loss with the sigmoid activation function.

\subsection{Hidden Layer Error}

For hidden layers, the error is propagated backward from the next layer:

\[
\delta_j = \left( \sum_k \delta_k w_{kj} \right)\sigma'(z_j)
\]

where:
\begin{itemize}
    \item $\delta_k$ is the error term from the next layer,
    \item $w_{kj}$ is the weight connecting neuron $j$ to neuron $k$,
    \item $\sigma'(z_j)$ is the derivative of the sigmoid function.
\end{itemize}

The sigmoid derivative is:

\[
\sigma'(z) = \sigma(z)(1-\sigma(z))
\]

\subsection{Gradient Computation}

The gradient of the loss with respect to each weight is computed as:

\[
\frac{\partial L}{\partial w_i} = \delta \times x_i
\]

The gradient with respect to the bias is:

\[
\frac{\partial L}{\partial b} = \delta
\]

\subsection{Weight Update}

After computing gradients for all training samples, the parameters are updated using gradient descent:

\[
w = w - \eta \frac{\partial L}{\partial w}
\]

\[
b = b - \eta \frac{\partial L}{\partial b}
\]

where $\eta$ is the learning rate.

\subsection{Training Process}

The training procedure consists of the following steps:

\begin{enumerate}
    \item Perform forward propagation to compute predictions.
    \item Compute the loss value.
    \item Perform backpropagation to calculate gradients.
    \item Update weights and biases using gradient descent.
    \item Repeat until the loss reaches a predefined threshold or the maximum number of epochs is reached.
\end{enumerate}

\section{Experiment with the XOR data set and the house price-size}

The implemented neural network uses the sigmoid activation function:

\[
\sigma(z) = \frac{1}{1 + e^{-z}}
\]

The sigmoid function produces outputs only within the range:

\[
0 < \sigma(z) < 1
\]

This activation function is suitable for binary classification problems such as the XOR dataset, where the target outputs are either 0 or 1.

However, the house price-size dataset is a regression problem, where the target values are continuous numerical values such as house prices. Since sigmoid restricts the output to the interval $(0,1)$, the network cannot naturally predict large continuous values.

For regression tasks, the output layer should typically use a linear activation function:

\[
\hat{y} = z
\]

which allows the network to predict any real-valued output.

Therefore, while the sigmoid activation function works well for the XOR dataset, it is not appropriate for house price prediction.
\end{document}